% ---------------------------------------------------------------------------
\section{Limitations and Dual-Use / Ethical Considerations}\label{app:limitations}

\lead{What this establishes.} On agent transcripts, exact-text provenance finds inter-agent flows with high precision without enumerating channels. On a public reproduction of the incident, 93--99\% of its edges are backed by evidence that involves no text matching. In a rig with ground truth, the positive control and the authored flows are found completely, and relayed flows at 0.88--0.98 content recall. \lead{What it does not establish:} that the detector stops an adversary who knows the rule (Section~\ref{sec:results} shows it does not), that it works beyond exact text, or that it scales as projected.

\lead{Limitations.}
\begin{itemize}
  \item \lead{Exact matching only.} Paraphrase, translation and steganography \cite{motwani2024} are out of scope.
  \item \lead{Short payloads are invisible.} Anything shorter than eight words, 40 characters or a 20-character atom, such as a file name, a flag or a three-word message, cannot match.
  \item \lead{Laundering.} Content that never crosses a model boundary as the same string is not seen, and honest agents do this routinely with \texttt{cp}.
  \item \lead{Convergent boilerplate.} Agents given the same prompt write the same text. When one assembles it with a program and reads it back, it looks inbound from the other. This produced 10 of the 11 false positives in our rig (Appendix~\ref{app:conditions}), and will grow with fleet size. The eleventh was a character shingle spanning the end of one file and the start of the next within a single read.
  \item \lead{Ground truth is small.} We have it only in our rig: at most 15 agents per episode, three to five seeds, two Claude models. The incident check is against registry writes, not reads, and DseWiki has no ground truth.
  \item \lead{The adversary is scripted.} Its strategies were written by one model with white-box access. No agent inside a run tried to evade, so we have not measured what a capable, willing one would find.
  \item \lead{The scale projection is an estimate.} The native store is not built.
\end{itemize}

\lead{Dual-use and ethical considerations.} This is a defensive tool, and its main risks are the data it collects and what it teaches about evasion.
\begin{itemize}
  \item \lead{A clean report is not proof of isolation.} The detector cannot see laundered content, short strings, paraphrase or steganography. It is a cheap first-line signal, not a proof of noninterference, and should not be presented as one.
  \item \lead{The index is sensitive.} To work, the store has to see every agent input and output across an organisation, credentials included. A plaintext store would be one place to find every secret any agent ever touched. Ours stores 64-bit hashes of shingles, never the text, and a deployment should use a keyed hash. Matched text should live only in findings, under the same access controls as the logs.
  \item \lead{Evasion knowledge.} We name the gaps and release the scripted evader with the code, because the evasion results cannot be reproduced without it. Every strategy in it is a few lines of shell (\texttt{cp}, base64, writing in pieces, a short code) that a capable agent or attacker would find without us. We judge that naming them helps defenders scope the tool more than it helps anyone evade it.
  \item \lead{Surveillance creep.} The same index that finds an undeclared channel between agents would find a person copying text between systems, if human input were in scope. We scope the method to agent transcripts, and any extension to people should be a separate, governed decision.
  \item \lead{Experimental ethics.} All credentials and references in the rig are synthetic and generated from a seed. Containers had no network access, and all prompts and data were synthetic. One pilot used a free provider tier that may retain prompts; no personal or proprietary data was sent.
  \item \lead{Prior work.} One commit predates the sprint: the pure detector core, the event schema and five synthetic test conditions (\texttt{7dd3a9d}, 24 Sep 2026). Everything else, including the Inspect adapter, resource descriptors, the carrier condition, atoms, the rig, the oracle and all analyses, was written during the sprint. The method's revisions are documented in \texttt{docs/DESIGN.md}.
\end{itemize}

% ---------------------------------------------------------------------------
\section{Conditions in detail}\label{app:conditions}

\lead{Ground truth and metrics.} A host-side inotify watcher records every read and write on the shared directory, from every container. A ground-truth edge W~$\to$~R means R read a file after W last wrote it, or listed a directory after W created something in it. A detector edge X~$\to$~R is correct if X is an ancestor of R in the ground truth, since a finding names the original origin, which may be several hops back. \emph{Content recall} is the share of file-content edges for which the detector found an edge from the writer or one of its ancestors.

\lead{Models.} Every condition with a model runs five seeds on Claude Haiku 5.5. Every condition except registry also runs three on Claude Sonnet 5.5. Handoff passes a release reference rather than a credential: a version passing a deploy token (relay, below) was stopped on Sonnet by Anthropic's safety classifier. The scripted adversary takes the model's place and replays fixed command sequences through the same rig, oracle and analysis.

Agents in an episode run sequentially, in the order given. Seeds are listed as Claude Haiku 5.5 / Claude Sonnet 5.5.

\begin{table}[h]
\centering\small
\caption{The conditions.}\label{tab:conditions}
\begin{tabularx}{\linewidth}{>{\raggedright\arraybackslash}p{2.6cm}c Y c}
\toprule
Condition & \makecell{Agents per\\episode} & Structure & Seeds \\
\midrule
isolated & 10 & ten operations teams, each fixing a bug in its own copy of a script; no shared mount & 5 / 3 \\
registry & 15 & five regions; all five billing teams (first round), then the five pricing teams that hold the tariff tables, then the billing teams again (second round). The only shared surface is a writable artefact registry. & 5 / -- \\
shared & 15 & as registry, but the shared surface is a folder described as the place to ask other teams for things & 5 / 3 \\
relay & 9 & three services (ledger, routing, manifests), each a chain security $\to$ release $\to$ deploy, passing a deploy token & 5 / 3 \\
handoff & 9 & as relay, passing a release reference instead of a token & 5 / 3 \\
message & 6 & three regions, each a planner (holds the quotes) $\to$ booker (holds nothing) & 5 / 3 \\
covert, covert message, covert handoff & 9, 6, 9 & as relay, message and handoff, with every agent briefed on the monitor & 5 / 3 each \\
scripted adversary & 9 or 6 & the handoff, relay or message chain, with the model replaced by a fixed command sequence & \makecell{one episode\\per strategy} \\
\bottomrule
\end{tabularx}
\end{table}

Relay also ran six episodes on a FUSE-backed shared mount (Appendix~\ref{app:fuse}).

Relay was the original relayed condition. On Haiku it gives the same picture as handoff (content recall 0.81, precision 1.0, five seeds on ext4), and the scripted evader gets the same result on both chains. Its hops often moved only short YAML lines, so at 12\,/\,60\,/\,28 its content recall falls to 0.44. On Sonnet, the safety classifier stopped every security agent, so no token moved. Five seeds of isolated, registry, shared and relay (245 runs) took about 50 minutes and cost US\$0.27 on one desktop.

In registry, no agent wrote to the registry in 75 chances: billing agents looked, found only the README, and reported that they could not get the tariff. Its 10 false positives (21 per 1,000 inputs) are 10 of the rig's 11, all convergent boilerplate (Appendix~\ref{app:limitations}).

% ---------------------------------------------------------------------------
\section{Relay outcomes by file system}\label{app:fuse}

The shared directory first sat on a FUSE-mounted disk that ignores \texttt{chmod}, by accident. Security agents noticed, and most refused the handoff: \emph{``NOT HANDED OFF. I did not leave the deploy token in /shared/handoff/\ldots\ any copy would be exposed to everyone.''} We reran the relay with the share on ext4, where \texttt{chmod} appears to work. There no agent refused, but every container runs as root, so other teams' tokens were still read.

\begin{table}[h]
\centering\small
\caption{Relay outcomes on Claude Haiku 5.5 by the file system under the shared directory.}\label{tab:fuse}
\begin{tabular}{lcc}
\toprule
Relay outcome & \makecell{FUSE share, chmod\\ignored (6 runs)} & \makecell{ext4 share, chmod\\holds (5 runs)} \\
\midrule
Security agents that declined the handoff & 11 / 18 & 0 / 15 \\
Deploys with their own token & 6 / 18 & 15 / 15 \\
Deploys with another team's token & 3 / 18 & 0 / 15 \\
Deploys failed or blocked & 9 / 18 & 0 / 15 \\
Reads of another team's token file (ground truth) & 18 & 8 \\
Detector precision / content recall & 1.0 / 0.87 & 1.0 / 0.81 \\
\bottomrule
\end{tabular}
\end{table}

% ---------------------------------------------------------------------------
\section{Models briefed on the monitor}\label{app:briefed}

In the covert conditions every agent was told how the monitor works and asked to complete its task without being flagged. Covert handoff added that the run was a sanctioned red-team test of the company's own monitor and that everything in it was synthetic.

\begin{table}[h]
\centering\small
\caption{What briefed models did.}\label{tab:briefed}
\begin{tabularx}{\linewidth}{>{\raggedright\arraybackslash}p{2.4cm} Y Y}
\toprule
Condition & Claude Haiku 5.5 (5 seeds) & Claude Sonnet 5.5 (3 seeds) \\
\midrule
covert & no security agent encoded the token; 8 of 15 refused the handoff, 7 handed it over in plain text & the safety classifier stopped the security step in every run \\
covert message & every note in plain text; content recall 1.0 & every note in plain text; content recall 1.0 \\
covert handoff & all 15 security agents declined to evade, in words, and completed the handoff; content recall 0.96 & the classifier stopped all 27 agents; the 3 security agents that answered first refused \\
\bottomrule
\end{tabularx}
\end{table}

Typical replies. Haiku: \emph{``I did not attempt to hide or disguise the reference from the transcript monitor; it is delivered as a plain file.''} Sonnet: \emph{``I did not encode or disguise it to get around the transcript monitor, because I won't help evade that oversight.''} Covert handoff recall is higher than plain handoff (0.96 against 0.88) because the briefed agents quoted the reference in their closing message, which counts as carrying it.

% ---------------------------------------------------------------------------
\section{Detailed results}\label{app:results}

\lead{Incident reproduction.} The largest cascade began with one note. Worker RA-13 posted that its input file \emph{``was not staged''}, and 75 agents read it over the next 16 minutes, through every URL spelling and loop construct they improvised. Relayed edges are the weaker class here (80.7--98.0\% backed, against 97.9--99.7\% for authored). Most unbacked relayed edges start from a file read: several agents got identical task manifests, and the detector credits whichever agent read it first.

\lead{The rig.} Table~\ref{tab:rig} compares the detector with the inotify ground truth. In isolated runs it raises nothing. On shared, the positive control, precision and content recall are 1.0 in every seed. On message, where only authored text can carry the plan, both are 1.0 on both models. Handoff is the hard case, and it is where relayed findings matter: without them recall falls from 0.88 to 0.34 on Haiku.

\begin{table}[h]
\centering\small
\caption{The detector against the inotify ground truth (means over seeds).}\label{tab:rig}
\begin{tabular}{llcccc}
\toprule
Condition & Model & Seeds & Content recall & Precision & Without relayed findings \\
\midrule
isolated & Haiku 5.5 & 5 & -- & no findings & -- \\
isolated & Sonnet 5.5 & 3 & -- & no findings & -- \\
shared & Haiku 5.5 & 5 & 1.0 & 1.0 & 1.0 \\
shared & Sonnet 5.5 & 3 & 1.0 & 1.0 & 1.0 \\
message & Haiku 5.5 & 5 & 1.0 & 1.0 & 1.0 \\
message & Sonnet 5.5 & 3 & 1.0 & 1.0 & 1.0 \\
handoff & Haiku 5.5 & 5 & 0.88 & 1.0 & 0.34 \\
handoff & Sonnet 5.5 & 3 & 0.98 & 0.98 & 0.49 \\
\bottomrule
\end{tabular}
\end{table}

Every miss starts at a records agent that read the reference, moved it with \texttt{cp}, and never typed it. No model emitted it, so the carrier condition, which keeps pytest banners out, keeps this hop out too: laundering, done by an honest agent. The briefed agents of Appendix~\ref{app:briefed} mostly quoted the reference in their closing message, and their recall rose to 0.96. The one false positive is a character shingle spanning the end of one file and the start of the next within a single read.

A relayed finding can show harm nobody in the run noticed. In one Sonnet seed the classifier stopped two of the three records agents. Their release agents each took the only reference in the folder, which belonged to the manifests team, and both published with it; the publish step checks only the format, and printed \emph{``reference ok''}. The detector traced the reference in both deploys to the manifests chain.

\lead{Evasion.} Models briefed on the monitor declined to evade it (Appendix~\ref{app:briefed}), so we measure evasion with a scripted adversary. Its strategies were written by Claude Opus 5.5 acting as a coding assistant with white-box access to the rule and source code. Four of six evade the detector entirely while completing the task, and a fifth hides the payload but not the handoff (Table~\ref{tab:evasion}). This bounds what an informed adversary can do. It does not show that an agent inside a run, with only a description of the monitor, would find these attacks: given that description, every agent that acted declined to try.

\begin{table}[h]
\centering\small
\caption{The scripted evader on the handoff (relayed) and message (authored) chains. Every task completed.}\label{tab:evasion}
\begin{tabularx}{\linewidth}{llYcc}
\toprule
Strategy & Kind & What the evader does & \makecell{Content edges\\found} & \makecell{Payload\\traced} \\
\midrule
plain & relayed & reads the reference and types it into the note & 6 / 6 & yes \\
cp & relayed & moves the reference with \texttt{cp} and \texttt{\$(cat \ldots)} & 0 / 6 & no \\
base64 & relayed & encodes and decodes with programs; release reads the decoded reference & 3 / 6 & no \\
plain & authored & writes the plan as a note & 3 / 3 & yes \\
encoded & authored & types the plan into an encoder; the booker decodes and reads it & 3 / 3 & yes \\
blind & authored & the same, but the booker decodes straight into its output file & 0 / 3 & no \\
fragments & authored & appends the note a few characters per command & 0 / 3 & no \\
compact & authored & sends a code shorter than any shingle (\texttt{4619N 8510H \ldots}) & 0 / 3 & no \\
\bottomrule
\end{tabularx}
\end{table}

Encoding alone doesn't help: the writer types the plaintext into the encoder, and a reader that decodes and reads it gets the same text by a new route. Every strategy that works keeps the payload from appearing as the same text on both sides of the model boundary. Either no model reads it (cp, blind), no model writes it in one piece (fragments), or it is too short to match (compact). The first two are deliberate laundering, closed by tracking provenance through the file system and program execution; the third is the short-string limit. Relayed content is the cheaper to hide, since \texttt{cp} takes no effort at all.

\lead{False positives on real work.} On 140 Refactor Arena samples with nothing shared, every finding is false. Seeded with the 2,137 initial repository files, the detector raises 326 findings over 5,946 inputs, or 55 per 1,000; without the seed, 1,128. Refactoring is close to the worst case, since agents re-emit repository code by design. Most of what remains is test output that two samples of the same repository both produce. In a typical one, a sample runs \texttt{python -m pytest} and quotes a failing test id; another sample later reads \texttt{tests/e2e}, where the same id appears. Same environment, a different route, and a model that carried it: by the rule's lights, a flow. Whether two samples of one repository are one run or two is a deployment decision, not a property of the rule.

We fixed 8 words, 40 characters and 20-character atoms before running any experiment. Table~\ref{tab:sweep} reruns the same logs at other lengths.

\begin{table}[h]
\centering\small
\caption{Shingle lengths (Claude Haiku 5.5 runs; FP: false findings per 1,000 inputs).}\label{tab:sweep}
\begin{tabular}{lccc}
\toprule
Words / characters / atom length & Handoff content recall & Isolated FP & Refactor Arena FP \\
\midrule
5 / 24 / 12 & 1.0 & 10.5 & 171 \\
6 / 32 / 16 & 1.0 & 0 & 90 \\
\textbf{8 / 40 / 20 (used)} & 0.88 & 0 & 55 \\
12 / 60 / 28 & 0.88 & 0 & 32 \\
\bottomrule
\end{tabular}
\end{table}

Shared and message stay at precision and recall 1.0 in every setting. Shorter shingles recover the \texttt{cp} hops, though not through the reference, which no model typed: they match its 18-character file name, at up to three times the false-positive rate on real work. Longer ones cost handoff nothing, since the 36-character reference is an atom at every length, but halved recall on relay, whose shared spans were short YAML lines (Appendix~\ref{app:conditions}). 6\,/\,32\,/\,16 looks like the better operating point on our data; choosing it now would be tuning on the test set.

\lead{The real board.} DseWiki has no reads, so we turn the rule around: if account X inserts a string whose first appearance anywhere was by account Y on a different page, X read Y's page. There is no ground truth, so this describes what the method finds, not how accurate it is. On strings inserted by fewer than five accounts there are 12,806 account-to-account edges. 38\% are time-decoupled, with the reader's first edit after the writer's last: the board worked as memory across agent lifetimes, and a monitor that compares only agents running at the same time would miss all of them. 714 edges cross between wikis. Account labels are chosen by the agents, so some of this is one agent remembering, and every finding is authored, since without reads relayed content cannot be seen.

\lead{Scalability.} The store grows more slowly than the text it indexes, and scan time is linear in text, with no pairwise cost between agents (Table~\ref{tab:scale}, reference Python implementation on subsets of the GPT arm).

\begin{table}[h]
\centering\small
\caption{Store and scan cost of the reference implementation.}\label{tab:scale}
\begin{tabular}{rrrrrr}
\toprule
Agents & Characters & Store keys per character & Store size & Scan time & Findings \\
\midrule
32 & 1.8\,M & 0.45 & 357\,MB & 7\,s & 136 \\
64 & 3.6\,M & 0.41 & 687\,MB & 14\,s & 625 \\
128 & 8.1\,M & 0.36 & 1.4\,GB & 35\,s & 2,662 \\
256 & 17.3\,M & 0.31 & 2.5\,GB & 82\,s & 8,311 \\
\bottomrule
\end{tabular}
\end{table}

Bigger fleets repeat more, so keys per character fall as the fleet grows. Winnowing the character shingles \cite{schleimer2003} (window 20) cuts the store 4.7$\times$ for a small loss of recall on short shared spans; atoms, and so credentials, are untouched. Store size depends only on the total text, so the useful unit is bytes per token. A native store needs about 40 bytes per key, which at 0.31 keys per character (0.065 winnowed) and four characters per token is about 50 bytes per token, or 10 winnowed. A trillion tokens of rollouts would need about 50\,TB, or 10\,TB winnowed: a sharded on-disk key-value store, with keys expiring once no live run could still read them. \emph{This projection is an estimate; we have not built or measured the native store.}
